\section{Squarefree local support}
\label{sec:squarefree-support}

For \(r\) with \(L_1(r)L_2(r)\ne0\), define
\begin{equation}
 \Qsf(r)
 =
 \mu^2(|L_1(r)|)\mu^2(|L_2(r)|).
\label{eq:squarefree-support}
\end{equation}
At an integral zero of either form, set \(\Qsf(r)=0\).  This
convention changes at most two values and is convenient for density
statements.

\subsection{Finite local support}

For \(z\ge2\), let
\[
 P(z)=\prod_{p\le z}p^2
\]
and define the truncated local indicator
\begin{equation}
 \Qsf_z(r)
 =
 \prod_{p\le z}
 \ind_{r\bmod p^2\notin\Rset_{p,2}}.
\label{eq:truncated-support}
\end{equation}

\begin{theorem}[Exact finite local density]
\label{thm:finite-local-density}
The function \(\Qsf_z\) is periodic modulo \(P(z)\), and
\begin{equation}
 \frac1{P(z)}
 \sum_{r\bmod P(z)}\Qsf_z(r)
 =
 \prod_{p\le z}
 \left(1-\frac{\nu_{p,2}}{p^2}\right).
\label{eq:finite-local-product}
\end{equation}
Every factor on the right is positive.
\end{theorem}

\begin{proof}
For each \(p\le z\), exactly \(p^2-\nu_{p,2}\) residue classes
modulo \(p^2\) are allowed.  The moduli \(p^2\) are pairwise
coprime, so the Chinese remainder theorem multiplies the allowed
counts.  Finally, \(\nu_{p,2}\le2<p^2\), including at \(p=2\).
\end{proof}

\begin{corollary}[No squarefree local annihilation]
\label{cor:no-squarefree-local-annihilation}
For every finite set of primes, there is a residue class satisfying
\(p^2\nmid L_i(r)\) for \(i=1,2\) and every prime in the set.
Quantitatively,
\begin{equation}
 \prod_{p\le z}
 \left(1-\frac{\nu_{p,2}}{p^2}\right)
 \ge
 \prod_{p\le z}\left(1-\frac2{p^2}\right)>0.
\label{eq:finite-universal-lower}
\end{equation}
\end{corollary}

\begin{proof}
Use \(\nu_{p,2}\le2\) in
\eqref{eq:finite-local-product}.
\end{proof}

\subsection{The Euler product and fixed-form density}

\begin{definition}[Squarefree singular factor]
\label{def:squarefree-singular-factor}
For a fixed primitive affine pair, put
\begin{equation}
 \Mloc(L_1,L_2)
 =
 \prod_p
 \left(1-\frac{\nu_{p,2}}{p^2}\right).
\label{eq:squarefree-singular-factor}
\end{equation}
\end{definition}

\begin{proposition}[Absolute convergence and uniform positivity]
\label{prop:euler-product-positive}
The product \eqref{eq:squarefree-singular-factor} converges to a
positive number and satisfies
\begin{equation}
 \boxed{
 1\ge\Mloc(L_1,L_2)
 \ge
 \prod_p\left(1-\frac2{p^2}\right)>0.}
\label{eq:universal-squarefree-lower}
\end{equation}
\end{proposition}

\begin{proof}
Since \(0\le\nu_{p,2}\le2\) and
\(\sum_p p^{-2}<\infty\), the logarithm of the Euler product
converges absolutely after finitely many initial factors.  Every
factor is positive.  The lower bound follows termwise.
\end{proof}

\begin{theorem}[Simultaneous squarefree density for a fixed pair]
\label{thm:fixed-pair-squarefree-density}
For every fixed primitive affine pair with \(h_0\ne0\),
\begin{equation}
 \lim_{R\to\infty}
 \frac1R\sum_{1\le r\le R}\Qsf(r)
 =
 \Mloc(L_1,L_2).
\label{eq:fixed-pair-density}
\end{equation}
\end{theorem}

\begin{proof}
Fix \(z\).  Periodicity and
\cref{thm:finite-local-density} give
\[
 \lim_{R\to\infty}
 \frac1R\sum_{r\le R}\Qsf_z(r)
 =
 \prod_{p\le z}
 \left(1-\frac{\nu_{p,2}}{p^2}\right).
\]
It remains to remove primes \(p>z\).

Let
\[
 B_R=\max_{1\le r\le R}
 \max\{|L_1(r)|,|L_2(r)|\}.
\]
For the fixed affine pair, \(B_R\ll_{L_1,L_2}R\).  Away from the at
most two integral zeros, a discrepancy
\(\Qsf_z(r)\ne\Qsf(r)\) requires
\[
 p^2\mid L_i(r)
\]
for some \(i\in\{1,2\}\) and prime
\(z<p\le\sqrt{B_R}\).  If \(p\) divides the relevant slope, there is
no root by \cref{lem:one-linear-congruence}; otherwise there is one
residue class modulo \(p^2\).  Hence the number of discrepancies is
at most
\[
 2R\sum_{p>z}\frac1{p^2}
 +2\pi(\sqrt{B_R})+O(1)
 \ll
 \frac{R}{z}+\sqrt R+O(1).
\]
After division by \(R\), first let \(R\to\infty\), then
\(z\to\infty\).  The finite products converge to
\eqref{eq:squarefree-singular-factor}, proving the theorem.
\end{proof}

\begin{remark}[What the density proves]
\label{rem:density-scope}
The theorem proves that a \emph{fixed} primitive pair has positive
simultaneous squarefree density.  It does not prove cancellation of
Mobius or Liouville signs.  It also does not give a uniform
translated-short-interval theorem for coefficients depending on
\(X\).
\end{remark}

\subsection{Why growing physical fibers need a separate theorem}

\begin{proposition}[Elementary tail bound on one finite interval]
\label{prop:finite-interval-tail}
Let \(I\subset\Z\) be an interval of \(N\) consecutive integers, and
put
\[
 B=\max_{r\in I}\max\{|L_1(r)|,|L_2(r)|\}.
\]
Apart from integral zeros of the forms,
\begin{equation}
 \#\{r\in I:\Qsf_z(r)\ne\Qsf(r)\}
 \ll
 \frac{N}{z}+\sqrt B.
\label{eq:finite-interval-tail}
\end{equation}
The implied constant is absolute.
\end{proposition}

\begin{proof}
Repeat the union bound in the proof of
\cref{thm:fixed-pair-squarefree-density}.  A single soluble class
modulo \(p^2\) meets \(I\) at most \(N/p^2+1\) times, and only primes
\(p\le\sqrt B\) can occur.
\end{proof}

\begin{remark}[A limitation, not a counterexample]
\label{rem:tail-bound-limitation}
If the slopes and interval origins grow, then \(B\) may be much
larger than \(N^2\), and \eqref{eq:finite-interval-tail} gives no
useful relative error.  This does not prove that squarefree mass is
small; it proves only that the fixed-form density argument is not a
uniform physical theorem.  Likewise, the exact periodic density in
\eqref{eq:finite-local-product} need not be visible inside an
interval much shorter than \(P(z)\).
\end{remark}
